\section{Introduction}
\label{sec:introduction}

First-passage percolation, introduced by Hammersley and
Welsh~\cite{HammersleyWelsh1965}, equips a lattice with a random metric
by assigning independent passage times to its edges. A basic question asks
whether this metric can contain a doubly infinite minimizing path.
Such a path, called a \emph{bigeodesic}, is a random analogue of a
complete minimizing geodesic in a noncompact metric space.
The question was attributed to Furstenberg by
Kesten~\cite[Remark 9.22]{Kesten}.
The planar no-bigeodesics conjecture asserts that bigeodesics are absent
for iid continuous edge-weight laws; see
\cite[Conjecture 1 and the discussion preceding it]{DamronHanson}.
We prove the conjecture under the minimum-of-four second-moment
assumption.

Let \(\mathcal E\) denote the unoriented nearest-neighbor edges of
\(\Z^2\). Give these edges independent weights \(t_e\) with common
Borel probability law \(G\) on \([0,\infty)\). For vertices \(x,y\),
let \(T(x,y)\) be the infimum of the costs of finite nearest-neighbor
walks from \(x\) to \(y\), with each edge traversal charged its weight.

\begin{theorem}\label{thm:main}
Suppose that \(G\) has no atoms and that, for independent
\(X_1,X_2,X_3,X_4\) with law \(G\),
\begin{equation}\label{eq:minimum-moment}
 \EE\!\left[\min\{X_1,X_2,X_3,X_4\}^{\,2}\right]<\infty.
\end{equation}
Then, with probability one, there is no sequence
\((v_i)_{i\in\Z}\) of pairwise distinct vertices with consecutive
vertices adjacent and
\[
 \sum_{k=i}^{j-1}t_{\{v_k,v_{k+1}\}}=T(v_i,v_j)
 \qquad\text{for every }i<j.
\]
\end{theorem}

The conclusion excludes all bigeodesics in one probability-one event.
In particular, an end direction may depend on the environment.
The law \(G\) may be singular continuous or have bounded support,
and \(\EE t_e\) may be infinite. We assume no differentiability,
strict convexity, or curvature of the limit shape.
Theorem~\ref{thm:main} resolves the planar no-bigeodesics conjecture
positively in the class \eqref{eq:minimum-moment}; it does not
assert the moment-free formulation.

\subsection*{Background}
The study of infinite geodesics and Busemann functions developed through
work of Newman~\cite{Newman1995}, Hoffman~\cite{Hoffman2005,Hoffman2008},
and the stationary constructions of Damron and
Hanson~\cite{DamronHanson2014}. Directional results connect geodesics
to the geometry of the limit shape. Licea and Newman proved coalescence
for each direction in a deterministic set of full Lebesgue measure and
excluded bigeodesics whose two end-directions are prescribed in that
set~\cite{LiceaNewman}; see also
\cite[Theorem~3.1 and Remark~2 after Theorem~1]{DamronHanson}.
This coalescence result requires no curvature assumption.
The complementary assertion that every ray has a direction and that
rays exist in every direction was established by Newman under an
exponential moment and uniform curvature
\cite[Theorem~2.1]{Newman1995}. Damron and Hanson subsequently obtained
one-end exclusions in fixed tangent-contact sectors under local
differentiability assumptions~\cite[Theorem~1]{DamronHanson}.
Their full iid hypothesis A1$'$ already allows
\eqref{eq:minimum-moment}; the distinction here is the absence of shape
assumptions and the simultaneous treatment of every possible path.

Wehr and Woo excluded bigeodesics in the half-plane restricted metric
under a finite individual-edge mean~\cite{WehrWoo1998}.
Auffinger, Damron, and Hanson removed that moment assumption for
continuous iid weights~\cite[Appendix~C]{AuffingerDamronHanson2015}
and proved coalescence and finite backward clusters for the associated
boundary-limit geodesic forest
\cite[Theorem~1.5]{AuffingerDamronHanson2015}.
Brito, Damron, and Hanson proved the absence of doubly infinite directed
paths in geodesic graphs obtained as limits of point-to-hyperplane
constructions in a fixed direction~\cite{BritoDamronHanson2023}.
Their iid hypothesis also gives \eqref{eq:minimum-moment} in two
dimensions, but their selected graphs need not contain every geodesic
of the metric. Alexander proved simultaneous absence of bigeodesics
under additional exponential-moment, uniform-curvature, and quantitative
fluctuation assumptions~\cite{Alexander2023}.

Ahlberg and Hoffman developed a structural theory of random
coalescing geodesics that assigns an asymptotically linear Busemann
function to every ray, without imposing shape regularity
\cite{AhlbergHoffman}.
Their results also give uniqueness, coalescence, and backward
finiteness for each fixed functional in a deterministic closed set.
These conclusions, stated in \Cref{thm:structure}, are the structural
inputs to our proof. We use them with the boundary comparison in
Appendix~\ref{sec:boundary-comparison}, which replaces a boundary-expectation
step without requiring integrability of restricted passage times.
Related developments include sublinear growth of planar geodesic
trees~\cite{AhlbergHansonHoffman} and quantitative coalescence and
midpoint estimates under stronger distributional and geometric
hypotheses~\cite{DembinElboimPeled}.

There is a quantifier issue that those inputs alone do not settle:
almost-sure uniqueness for each fixed functional does not imply
simultaneous uniqueness at every environment-dependent functional.
A bigeodesic, if one existed, could select such an exceptional
functional. Our argument treats these random functionals directly
and then uses an increasing passage-time cost to exclude all
remaining candidates.

\subsection*{The argument}
A ray \(p\) has a Busemann function
\[
 B_p(x,y)=\lim_{j\to\infty}
          \bigl(T(x,p_j)-T(y,p_j)\bigr)
\]
with a linear asymptote \(\rho_p\), which we call its label.
The labels of the two ends of a bigeodesic are opposite.
After a countable reduction, it suffices to exclude bigeodesics whose
upper labels lie in one compact sector of the dual unit circle.
The horizontal coordinate \(r(\rho)=\rho(e_1)\) orders this sector;
write \(W\) for its width in that coordinate.

The contradiction is a finite budget for passage-time increases.
On a horizontal grid of spacing \(l\), we mark sites by selecting
bigeodesics in finitely increased environments. A marked path acquires
an extra cost of at least \(M\) on an initial portion.
Partition the label sector into bins with deterministic endpoints.
The selected arms and the base rays used for comparison are required
to satisfy common deterministic cone bounds.
When the marks in each bin are sufficiently separated, comparison
with rays at those endpoints charges the increase to the bin's width.
Busemann increments telescope along the line, and the resulting bound is
\[
 Mq\le lW,
\]
where \(q\) is the common probability that a site is marked
(\Cref{prop:label-budget}). We construct such marks with \(q\) bounded
away from zero and \(l\) fixed while \(M\) becomes arbitrarily large.
Two issues require the rest of the proof: obtaining large increases
without losing the marks, and separating their labels without increasing
the grid spacing.

For the first issue, we raise finitely many edge weights in a widening
region around each site. Our low-to-high replacement is an instance of
the bounded-density Bernoulli perturbations of Bates and Chatterjee
\cite[Lemmas~2.3--2.4]{BatesChatterjee2020}.
Quadratic change-of-law bounds for spatially varying passage-time
perturbations also appear in
\cite[Section~2.3]{DembinElboimPeled}.
All densities here are relative to the original edge law, including
when it is singular. At height \(h\), we use the replacement scale
\[
 \delta_h=\frac{\eta}{(h+2)\log(h+2)}
\]
with one sufficiently small \(\eta>0\).
There are order \(h\) possible edges at that height. Thus the convergent
sum of \(h\delta_h^2\) keeps the change of law uniformly small, while
the divergent sum of \(\delta_j\) permits an arbitrarily large gain
along a path. The path must be chosen from the increased environment,
not in advance. Conditioning on that individual output gives an explicit product
law for the hidden raises along the selected path.

The second issue rests on a simultaneous statement about random labels:
two bigeodesics with the same label have the same pair of Busemann
functions (\Cref{cor:same-busemann}). To prove it, we first use planar
ordering and fixed bounding rays to confine every ray in the sector.
The labels admitting branching or distinct Busemann functions form a
countable set in each environment. Resampling complementary half-planes
independently, we transfer rays across a distant cut while preserving
their exact labels and finite witnesses to exceptionality
(\Cref{lem:cut-transfer}). Countability and independence imply
coalescence at the opposite end of every exceptional label.
A one-edge version of the budget and a finite-tree count then exclude forks
(\Cref{prop:edge-count,prop:no-forks}). These steps give the
common Busemann pair without intersecting uncountably many fixed-label
probability-one events.

To use this conclusion, all sites share the same raised samples.
Passing from two individual perturbation regions to their union leaves
both selected paths and their costs unchanged. Choosing \(l\) large
enough once ensures that a strictly raised edge within a fixed near
radius of one site lies outside the other site's region. If their labels
agreed, their common Busemann pair would produce
a detour through that edge with arbitrarily small excess cost in the
union metric. Returning to the other individual metric would make this
detour strictly cheaper than its selected geodesic. Hence distinct
marked sites have distinct labels (\Cref{prop:distinct-marks}).

For each fixed gain \(M\), refining the label partition therefore
separates every mark from the finitely many neighbors within the
deterministic separation required by the budget comparison. The probability lost in
this separation tends to zero as the partition is refined. Only after
this refinement do we increase \(M\). The grid spacing remains fixed,
and \(Mq\le lW\) is impossible. No quantitative separation of the
random labels, and no rate of coalescence, is required.
